\documentclass[10pt,twocolumn]{article}

\usepackage[margin=0.72in,columnsep=0.24in]{geometry}
\usepackage{amsmath,amssymb,amsthm,booktabs,mathtools,microtype}
\usepackage[hidelinks]{hyperref}
\setlength{\emergencystretch}{1em}

\newcommand{\E}{\mathbb E}
\newcommand{\R}{\mathbb R}
\newcommand{\KL}{\mathrm{KL}}
\newcommand{\1}{\mathbf 1}
\newtheorem{proposition}{Proposition}
\newtheorem{corollary}[proposition]{Corollary}
\theoremstyle{remark}
\newtheorem{remark}[proposition]{Remark}

\title{Ternary Relation Fields from Protein Families\\
to Sequence Transformers}
\author{PRL/PRX working draft}
\date{July 2026}

\begin{document}
\maketitle

\begin{abstract}
Multiple-sequence alignments sample a high-dimensional categorical distribution shaped by
conservation, phylogeny, and epistatic constraints.  Standard attention routes information with a
nonnegative probability distribution over source positions, whereas graphical sequence models
assign signed, typed interactions to position pairs.  We study a ternary relation field in
which every candidate edge occupies one of three states, $-1$, $0$, or $+1$, representing negative
polarity, inactivity, or positive polarity relative to a fixed edge orientation.  A three-state
softmax supplies differentiable state probabilities; an amino-acid softmax separately normalizes
the reconstructed residue distribution.  The model decomposes each relation into activity,
conditional polarity, magnitude, and a reference-centered categorical operator.  We prove that the
soft ternary state is exactly equivalent to an activity--polarity factorization and identify when
its sign is redundant or meaningful.  A continuous signed energy already spans attraction, zero,
and repulsion, so ternarity does not add a new equilibrium interaction value.  Its possible content
is instead an exact support state, uncertainty over an oriented mechanism, and a quantized
computational control.  For associative retrieval we derive a constrained three-state potential
whose mean gate is the derivative of a smooth thresholded energy; the tied-memory update follows
from a concave--convex procedure and monotonically decreases that energy.  We also distinguish the
common mean-gate relaxation from exact
marginalization over latent energy states.  We study two realizations of the same state space: a
static family-level graph fitted to an MSA and an ESM-style dynamic mixer whose edge states depend
on a single sequence.  An optional teacher--student bridge transfers family relations during
training without requiring an MSA at inference.  We formulate the phase diagram and cross-family
controls needed to determine whether ternary relations improve likelihood, epistasis recovery,
transfer, and realized compute at matched model budgets.
\end{abstract}

\section{Introduction}

A protein family is an evolutionary ensemble of categorical sequences constrained by fold,
function, and historical contingency.  A multiple-sequence alignment (MSA) provides a finite and
phylogenetically correlated sample
\begin{equation}
X^{(n)}=(X_1^{(n)},\ldots,X_L^{(n)}),
\qquad X_i^{(n)}\in\{1,\ldots,q\},
\label{eq:msa}
\end{equation}
from that ensemble.  Reconstructing its distribution is useful for two different reasons.  As a
predictive problem, it asks whether masked or corrupted residues can be recovered.  As an inverse
problem, it asks which one-site and multi-site constraints are supported by the observations.  The
second question is harder: conservation, phylogeny, finite sampling, and epistasis can all improve
prediction while implying different interaction structures.

Protein sequence models make different compromises between these two objectives.  A
position-specific model represents conservation but not epistasis.  Potts and related inverse
statistical models attach a typed $q\times q$ coupling table to each position pair and have been
successful in contact and mutational analyses \cite{morcos2011,ekeberg2013}.  Protein language
models instead optimize masked-token prediction with deep, input-dependent representations
\cite{rives2021}; MSA Transformer extends this strategy to alignment tensors \cite{rao2021}.
Dauparas et al. showed that PSSMs, Potts pseudolikelihood, Gaussian models, and autoencoders can be
placed in one reconstruction framework \cite{dauparas2019}.  These connections clarify the loss
being optimized, but they do not make every internal edge weight an identified interaction.

For protein reconstruction, the unresolved issue is the state assigned to a candidate relation.
A Potts table is a signed,
typed interaction field, but its support is normally imposed by regularization or post hoc
thresholding.  Standard attention learns dynamic routing, but its source-softmax distributes one
unit of nonnegative mass across positions and generally has no exact inactive edge
\cite{vaswani2017}.  Negative effects remain possible through value and output projections, yet
their polarity is implicit and is contracted immediately into node states.  Thus relation
existence, signed routing, categorical response, and predictive uncertainty are present in the
models, but are not exposed as separate, jointly auditable variables.

We ask whether the relation state itself should be categorical.  For every candidate edge we use a
three-state distribution over negative polarity, inactivity, and positive polarity.  This edge
distribution is distinct from the categorical output distribution over amino acids:
\begin{equation}
\boxed{
\begin{gathered}
\text{edge-state softmax over }\{-1,0,+1\}\\
\neq\\[-2pt]
\text{residue softmax over }\{1,\ldots,q\}.
\end{gathered}}
\label{eq:two-softmaxes}
\end{equation}
The first softmax represents uncertainty about one relation; it is normalized across three states
for a fixed edge and does not force different source positions to compete.  The second softmax is
required to define a normalized conditional distribution for MSA reconstruction.

We study this construction at two scales.  At family level, one static ternary graph and one set of
reference-centered categorical fields are fitted to all rows of an MSA.  This setting permits direct
tests of support, polarity, and operator recovery under controlled sampling and phylogenetic
noise.  At sequence level, an ESM-style mixer predicts layer- and input-dependent ternary states
from hidden representations.  A leave-query-out teacher can transfer the family graph during
training, while inference remains single-sequence.  The two branches address complementary
questions: whether a ternary relation is statistically recoverable in an ensemble, and whether the
same organization is transferable as a dynamic computation.

Ternarity does not acquire physical or algorithmic content merely from the labels $-1$, $0$, and
$+1$.  A continuous signed coupling already spans negative, zero, and positive values.  If a full
categorical edge field is closed under negation, its scalar sign can be absorbed into the field.
Moreover, a soft mean gate is generally not the exact marginal of three latent energy models, and
an inactive state saves computation only when the hardened implementation skips the edge.  These
facts make binary-support signed energy, continuous signed routing, and ordinary sparse attention
decisive null models rather than weak ablations.

The individual ingredients also have direct predecessors.  Modern Hopfield networks express
softmax attention as associative retrieval from a log-sum-exp energy \cite{ramsauer2021}.  Neural Relational Inference learns
categorical latent edge types and may reserve one type for a non-edge \cite{kipf2018}.  Differential
Transformer subtracts two source-softmax maps to obtain continuous signed attention \cite{ye2024},
while SignGT learns signed graph attention values \cite{chen2026}.  Ternary Weight Networks
quantize parameters to $-1$, $0$, and $+1$ \cite{li2016}.  Table~\ref{tab:prior-art-boundary}
therefore treats the proposed model as a constrained synthesis, not as the invention of any one of
these primitives.

\begin{table*}[t]
\centering
\small
\begin{tabular}{p{0.18\textwidth}p{0.17\textwidth}p{0.17\textwidth}p{0.17\textwidth}p{0.21\textwidth}}
\toprule
Method & Discrete zero state & Signed relation & Protein MSA objective & Missing relative to the present test \\
\midrule
Potts/DCA \cite{morcos2011,ekeberg2013} & Only with an added sparsity model & Typed signed coupling
table & Family pseudolikelihood & No learned three-state uncertainty or transferable ESM router \\
NRI \cite{kipf2018} & Optional no-edge type & Multiple learned edge types & No & No categorical MSA
gauge, phylogeny audit, or protein transfer \\
Ternary Weight Networks \cite{li2016} & Zero parameter value & Signed quantized parameter & No &
Quantizes weights rather than inferring relation states \\
Differential Transformer \cite{ye2024} & No categorical inactive state & Difference of two continuous
attention maps & No & No per-edge state simplex or support-recovery target \\
MSA Transformer \cite{rao2021} & No & Implicit through value paths & Alignment MLM & No explicit
ternary relation recovery or hard sparse execution \\
Modern Hopfield attention \cite{ramsauer2021} & No exact inactive memory & Nonnegative memory
weights & Associative retrieval & No signed state or thresholded zero-retrieval regime \\
This work & Explicit $-1/0/+1$ edge state & Oriented polarity plus typed field & Family
pseudolikelihood and ESM MLM & Must demonstrate a new recovery law or resource frontier \\
\bottomrule
\end{tabular}
\caption{Prior-art boundary.  Every primitive of the architecture has a close predecessor.  The
scientific claim must therefore rest on recovery, transfer, or scaling behavior of their combined
protein-sequence realization.}
\label{tab:prior-art-boundary}
\end{table*}

The central hypothesis is consequently about collective behavior, not the three-way softmax.  The
ternary description is scientifically substantive only if it produces at least one phenomenon not
explained by its simpler nulls: a shifted recovery boundary relative to binary-support signed
energy, a cross-family scaling law for identifiability of the zero state, or transfer of calibrated
family relations into a single-sequence ESM with a measured accuracy--compute advantage.  The
associative realization should therefore be read as a signed, factorized memory: activity selects
whether a memory participates, while conditional polarity selects whether it is added or
subtracted.  The analysis below first separates probability, energy, activity, polarity, and typed edge value; it
then defines the two realizations and the falsification tests needed to distinguish these outcomes
from reparameterization.

\section{Probability, energy, and the role of ternarity}

Every strictly positive conditional distribution admits an energy coordinate
\begin{equation}
E_i(a\mid x_{-i})
=-\beta^{-1}\log P(X_i=a\mid x_{-i})+C(x_{-i}),
\label{eq:conditional-energy}
\end{equation}
where the additive function $C$ is unobservable.  This identity does not by itself make the score a
physical Hamiltonian.  A global equilibrium interpretation additionally requires that all
conditionals be compatible with one joint law, that undirected pair fields obey the relevant
transpose symmetry, and that the joint distribution have one partition function.  Ordinary
Transformer attention probabilities are routing coefficients, not configuration probabilities;
calling their negative logarithms energies is a coordinate choice rather than a mechanism.

For a fixed oriented scalar interaction mode $\Phi_{ij}$, a continuous coefficient
$\kappa_{ij}\in\R$ already represents the full progression
\begin{align}
\kappa_{ij}<0&\quad\text{(negative)},\nonumber\\
\kappa_{ij}=0&\quad\text{(inactive)},\nonumber\\
\kappa_{ij}>0&\quad\text{(positive)}.
\label{eq:continuous-three-regimes}
\end{align}

\begin{proposition}[No new energy values from ternary sign]
Let an interaction be $\kappa_{ij}\Phi_{ij}$ with unrestricted
$\kappa_{ij}\in\R$.  The parameterization
\begin{equation}
\kappa_{ij}=z_{ij}\rho_{ij},
\qquad z_{ij}\in\{-1,0,+1\},\quad \rho_{ij}\geq0,
\label{eq:sign-magnitude-reparameterization}
\end{equation}
represents exactly the same scalar energy family.
\end{proposition}

\begin{proof}
Given $(z,\rho)$, their product is real.  Conversely, take $z=\operatorname{sign}(\kappa)$ and
$\rho=|\kappa|$, with $z=0$ at $\kappa=0$.
\end{proof}

Thus a ternary edge is not a new attraction--repulsion physics when the continuous energy is
already signed.  It becomes substantive only if the zero label is used as an exact support and
execution decision, if $p^-,p^0,p^+$ quantify uncertainty over declared mechanisms, or if the
interaction is deliberately quantized.  For a Potts table the limitation is stronger: one edge
$G_{ij}(a,b)$ can favor some amino-acid pairs and disfavor others, so an intrinsic scalar label
``attractive'' or ``repulsive'' generally does not exist without an orientation functional.

\subsection{APC as an endpoint-activity null}

Average product correction (APC) removes a separable background from a nonnegative pair score
$S_{ij}$, such as a mutual-information score or a gauge-fixed coupling norm
\cite{dunn2008}.  In its standard form,
\begin{equation}
S_{ij}^{\mathrm{APC}}
=S_{ij}-\frac{S_{i\boldsymbol\cdot}S_{\boldsymbol\cdot j}}
{S_{\boldsymbol\cdot\boldsymbol\cdot}},
\label{eq:apc}
\end{equation}
where the dotted quantities are the corresponding row, column, and global means.  A useful null
model is
\begin{equation}
S_{ij}=a_i a_j+R_{ij},
\label{eq:endpoint-score-null}
\end{equation}
with $a_i$ the generic propensity of position $i$ to participate in a strong inferred relation and
$R_{ij}$ the pair-specific residual.  When the rank-one term dominates the means and the residual
has small row and column averages, Eq.~\ref{eq:apc} approximately removes $a_i a_j$ and retains
$R_{ij}$.

The ternary parameterization makes this interpretation testable.  Suppose the log-odds that an
edge is active decompose as
\begin{equation}
\log\frac{m_{ij}}{1-m_{ij}}
=u_i+u_j+\Delta_{ij},
\label{eq:activity-null}
\end{equation}
where $u_i$ and $u_j$ are endpoint propensities and $\Delta_{ij}$ is pair specific.  In the sparse
regime $m_{ij}\ll1$,
\begin{equation}
m_{ij}\simeq e^{u_i}e^{u_j}e^{\Delta_{ij}}.
\label{eq:sparse-activity-factorization}
\end{equation}
Under the endpoint-only null $\Delta_{ij}=0$, the active mass is therefore approximately rank one.
This yields the falsifiable interpretation
\begin{equation}
\boxed{
\begin{gathered}
\text{APC is a first-order estimator}\\
\text{of an endpoint-driven activity null.}
\end{gathered}}
\label{eq:apc-interpretation}
\end{equation}
It is not an identity, and it need not hold outside the sparse or approximately separable regime.
Endpoint propensity can collect conservation, entropy, effective alphabet size, finite MSA depth,
phylogeny, susceptibility, or neural query/key scale.  These causes must be separated by matched
null ensembles rather than assigned a physical meaning from APC alone.

This interpretation also explains why signed interactions can still show a large APC effect.
Contact pipelines usually compress $G_{ij}(a,b)$ to $\|G_{ij}\|$ or another nonnegative magnitude,
discarding polarity while retaining activity and endpoint susceptibility.  Categorical gauge
fixing and APC remove different nuisances: Eq.~\ref{eq:gauge} removes additive one-body components
inside each amino-acid coupling table, whereas APC removes an across-position separable mode after
that table has been compressed to a scalar.  The ternary state adds a diagnostic separation:
$m_{ij}$ measures activity, while $g_{ij}$ measures signed mean.  The proposed mechanism predicts
that the leading separable mode is stronger in $m_{ij}$ and in interaction magnitudes than in a
consistently oriented signed matrix $g_{ij}$.

\section{Ternary relation model}

For every candidate directed edge $j\to i$, let
\begin{equation}
(p_{ij}^{-},p_{ij}^{0},p_{ij}^{+})
=\operatorname{softmax}
\left((r_{ij}^{-},r_{ij}^{0},r_{ij}^{+})/\tau\right).
\label{eq:state-softmax}
\end{equation}
Define its active mass, signed mean, and polarity conditional on activity by
\begin{equation}
m_{ij}=p_{ij}^{+}+p_{ij}^{-},\qquad
g_{ij}=p_{ij}^{+}-p_{ij}^{-},\qquad
s_{ij}=\frac{g_{ij}}{m_{ij}+\varepsilon}.
\label{eq:activity-polarity}
\end{equation}
In a hard forward pass,
\begin{equation}
z_{ij}=\arg\max_{z\in\{-1,0,+1\}}p_{ij}^{z}
\in\{-1,0,+1\}.
\label{eq:hard-state}
\end{equation}
A Gumbel--softmax or straight-through estimator may train the discrete state
\cite{jang2017,maddison2017}; it is an optimization device rather than an exact-gradient claim.

The value carried by an edge must still distinguish amino-acid types.  Let $u_i(a)>0$ be a
reference distribution at position $i$.  A categorical field $G_{ij}(a,b)$ is placed in the
reference gauge
\begin{equation}
\sum_a u_i(a)G_{ij}(a,b)=0,
\qquad
\sum_b u_j(b)G_{ij}(a,b)=0.
\label{eq:gauge}
\end{equation}
Write $G_{ij}=\rho_{ij}\widehat G_{ij}$ with $\rho_{ij}\geq0$ and
$\|\widehat G_{ij}\|_{u_i,u_j}=1$.  The soft conditional reconstruction model is
\begin{align}
\ell_i(a)
&=h_i(a)+\sum_{j\ne i}g_{ij}\rho_{ij}
\widehat G_{ij}(a,X_j),
\label{eq:conditional-logit}\\
P_\theta(X_i=a\mid X_{-i})
&=\frac{\exp[\ell_i(a)]}
{\sum_{c=1}^{q}\exp[\ell_i(c)]}.
\label{eq:conditional-model}
\end{align}
The hard model replaces $g_{ij}$ by $z_{ij}$ and evaluates only active edges.  Equation
\ref{eq:conditional-model} is a pseudolikelihood conditional; a globally normalized joint law
requires compatible symmetric fields and a separate partition function.

The reconstruction likelihood in Eq.~\ref{eq:conditional-model} observes only the effective
coefficient
\begin{equation}
\kappa_{ij}^{\mathrm{eff}}=g_{ij}\rho_{ij}.
\label{eq:effective-coefficient}
\end{equation}
Consequently, soft activity, polarity, and magnitude are not separately identified by conditional
likelihood alone.  Their separation is imposed by the ternary prior, activity penalty, orientation,
hard-state calibration, or independent supervision.  Any empirical claim about recovered support
or polarity must therefore be evaluated against ground truth or held-out interventions rather than
read directly from a fitted value of $g_{ij}$.

\subsection{Mean gate is not latent-energy marginalization}

For one edge write $\Phi_z(a)=z\rho\widehat G(a,x_j)$ and let $Z_z$ normalize across candidate
residues.  If $z$ is a genuine latent physical regime, exact marginalization gives
\begin{equation}
P_{\mathrm{mix}}(a)
=\sum_{z\in\{-1,0,+1\}}p^z
\frac{\exp[h(a)+\Phi_z(a)]}{Z_z}.
\label{eq:latent-energy-mixture}
\end{equation}
The soft model used in Eq.~\ref{eq:conditional-model} instead computes
\begin{equation}
P_{\mathrm{mean}}(a)
=\frac{\exp[h(a)+\E[z]\rho\widehat G(a,x_j)]}
{Z_{\E[z]}}.
\label{eq:mean-gate-energy}
\end{equation}
In general $P_{\mathrm{mix}}\ne P_{\mathrm{mean}}$ because exponentiation and state-dependent
normalization do not commute with expectation.  Equation~\ref{eq:mean-gate-energy} is therefore a
differentiable mean-gate or mean-field relaxation.  A claim that $z$ is a fluctuating physical
state requires Eq.~\ref{eq:latent-energy-mixture}, hard categorical sampling, or an explicitly
justified approximation audit.

\subsection{Exact activity--polarity factorization}

\begin{proposition}[Ternary simplex coordinates]
Every distribution $(p^-,p^0,p^+)$ on $\{-1,0,+1\}$ has the representation
\begin{equation}
p^0=1-m,\qquad
p^+=\frac{m(1+s)}{2},\qquad
p^-=\frac{m(1-s)}{2},
\label{eq:simplex-factorization}
\end{equation}
with $m\in[0,1]$ and $s\in[-1,1]$.  It is unique for $m>0$, and the expected state is
$\E[z]=ms$.
\end{proposition}

\begin{proof}
Take $m=p^++p^-$ and $s=(p^+-p^-)/m$ for $m>0$.  The inverse formulas follow by solving their sum
and difference.  When $m=0$, the edge is inactive and polarity is unidentifiable but irrelevant.
\end{proof}

This factorization prevents a common error: $g_{ij}\approx0$ can mean either a confidently absent
edge ($m_{ij}\approx0$) or an active but polarity-uncertain edge
($m_{ij}>0$, $s_{ij}\approx0$).  Sparsity must penalize $m_{ij}$, not $|g_{ij}|$.

\subsection{A ternary associative potential}

The unconstrained three-logit model is useful for relation inference but does not by itself define
a state dynamics.  A more restrictive parameterization does.  Let $x_j\in\R^d$ be a stored
memory, $\xi\in\R^d$ the current associative state, and $u_j=x_j^{\mathsf T}\xi$.  For threshold
$\lambda_j\geq0$, constrain the three logits to
\begin{equation}
r_j^+=\beta(u_j-\lambda_j),\qquad
r_j^0=0,\qquad
r_j^-=\beta(-u_j-\lambda_j).
\label{eq:associative-ternary-logits}
\end{equation}
Their partition function and signed mean are
\begin{align}
Z_j(u)&=1+2e^{-\beta\lambda_j}\cosh(\beta u),\nonumber\\
g_{\lambda_j}(u)&=p_j^+(u)-p_j^-(u)
=\frac{2e^{-\beta\lambda_j}\sinh(\beta u)}
{1+2e^{-\beta\lambda_j}\cosh(\beta u)}.
\label{eq:ternary-associative-gate}
\end{align}
Thus the zero state is controlled by an explicit similarity threshold rather than by a third
arbitrary score.  Define the convex potential
\begin{equation}
\phi_{\lambda}(u)=\frac{1}{\beta}
\log\!\left[1+2e^{-\beta\lambda}\cosh(\beta u)\right],
\qquad \phi_{\lambda}'(u)=g_{\lambda}(u).
\label{eq:ternary-associative-potential}
\end{equation}
In the hard limit $\beta\to\infty$ away from the thresholds,
\begin{equation}
g_{\lambda}(u)\longrightarrow
\begin{cases}
+1,&u>\lambda,\\
0,&|u|<\lambda,\\
-1,&u<-\lambda,
\end{cases}
\label{eq:ternary-hard-threshold}
\end{equation}
so the construction retrieves aligned memories, suppresses weak matches, and subtracts
anti-aligned memories.

\begin{proposition}[Energy descent for tied ternary memories]
\label{prop:ternary-hopfield-descent}
For fixed memories $\{x_j\}_{j=1}^{N}$, fixed thresholds, and $\gamma>0$, let
\begin{equation}
E_{\mathrm{ter}}(\xi)=\frac12\|\xi\|_2^2
-\frac{\gamma}{N}\sum_{j=1}^{N}
\phi_{\lambda_j}(x_j^{\mathsf T}\xi).
\label{eq:ternary-hopfield-energy}
\end{equation}
The update
\begin{equation}
\xi^{t+1}=\frac{\gamma}{N}\sum_{j=1}^{N}
g_{\lambda_j}(x_j^{\mathsf T}\xi^t)x_j
\label{eq:ternary-hopfield-update}
\end{equation}
is a concave--convex procedure step and satisfies
$E_{\mathrm{ter}}(\xi^{t+1})\leq E_{\mathrm{ter}}(\xi^t)$.  Moreover,
$E_{\mathrm{ter}}$ is bounded below and coercive.
\end{proposition}

\begin{proof}
Each $\phi_{\lambda}$ is a log-sum-exp of affine functions and is therefore convex.  Hence
Eq.~\ref{eq:ternary-hopfield-energy} is a difference of convex functions.  Linearizing the convex
sum at $\xi^t$ gives a quadratic majorizer whose unique minimizer is
Eq.~\ref{eq:ternary-hopfield-update}; the standard majorization inequality proves monotone descent.
Finally, $\phi_{\lambda}(x_j^{\mathsf T}\xi)$ grows at most linearly in $\|\xi\|_2$, whereas the
positive quadratic term grows quadratically, proving coercivity and a finite lower bound.
\end{proof}

This energy is a Lyapunov function for a hidden-state retrieval dynamics.  It is not the Potts or
pseudolikelihood energy of an amino-acid sequence, and $-\log p_j^z$ is not a third equivalent
energy.  Exponential-capacity results for modern Hopfield networks also cannot be transferred
unchanged to correlated MSA rows; protein memories violate the random-pattern assumptions that
make those bounds possible \cite{ramsauer2021}.

\subsection{When the ternary sign is meaningful}

\begin{proposition}[Sign redundancy]
If every active edge carries an unrestricted field $G_{ij}$ from a class closed under negation,
then replacing binary support by $z_{ij}\in\{-1,0,+1\}$ does not enlarge the represented function
class: the nonzero sign can be absorbed into $\widetilde G_{ij}=z_{ij}G_{ij}$.
\end{proposition}

\begin{proof}
For $z_{ij}=0$, delete the edge.  For $z_{ij}=\pm1$, closure under negation keeps
$\widetilde G_{ij}$ in the original field class and leaves the effective interaction unchanged.
\end{proof}

Thus a physical or biological interpretation of $+/-$ requires an orientation convention.  One
option is to select an anchor functional $\ell_{ij}$ and enforce
\begin{equation}
\ell_{ij}(\widehat G_{ij})>0.
\label{eq:orientation}
\end{equation}
The ternary sign then denotes polarity relative to the declared anchor.  Possible anchors include
a detached MSA teacher field, a fixed biochemical contrast, or a reproducibly chosen leading
singular direction.  Edges with vanishing or unstable anchors must be reported as unoriented.  A
scalar sign is not an intrinsic excitatory/inhibitory label for an unrestricted $q\times q$
categorical interaction.

\section{Reconstructing the MSA}

Let $w_n\geq0$ be phylogenetic sequence weights with $\sum_nw_n=1$.  A masked or
pseudolikelihood reconstruction objective is
\begin{equation}
\mathcal L_{\mathrm{rec}}
=-\sum_{n=1}^{N}w_n\sum_{i=1}^{L}
\log P_\theta(X_i^{(n)}\mid X_{-i}^{(n)}).
\label{eq:reconstruction-loss}
\end{equation}
To make the zero state computationally and statistically substantive, use
\begin{equation}
\mathcal L
=\mathcal L_{\mathrm{rec}}
+\lambda_0\sum_{i\ne j}m_{ij}
+\lambda_{\mathrm{sym}}\mathcal R_{\mathrm{sym}}
+\lambda_{\mathrm{cal}}\mathcal R_{\mathrm{cal}}.
\label{eq:total-loss}
\end{equation}
Here $\mathcal R_{\mathrm{sym}}$ enforces any joint-model reciprocity convention, while
$\mathcal R_{\mathrm{cal}}$ controls the soft-to-hard gap.  The expected active degree is
\begin{equation}
\bar k=\frac{1}{L}\sum_i\sum_{j\ne i}m_{ij},
\label{eq:expected-degree}
\end{equation}
and the hard realized degree is measured after Eq.~\ref{eq:hard-state}.  A low expected degree is
not a compute result unless the implementation actually skips zero edges.

Three reconstruction targets should be kept distinct:
\begin{enumerate}
\item conditional likelihood of held-out MSA residues;
\item recovery of known synthetic support, polarity, and categorical fields;
\item reproduction of ensemble observables such as marginals, connected correlations, contacts,
and mutational epistasis.
\end{enumerate}
Good masked likelihood alone does not establish interaction recovery.  Conversely, a field that
recovers synthetic couplings may fail under finite-depth, phylogenetically correlated MSAs.

\section{Two realizations of the ternary field}

The same three-state variable can be used at two statistical scales.  They share state semantics
and ablations, but they are not the same estimator.

\subsection{Static family-level graph}

In the family model, $r_{ij}^{z}$, $\rho_{ij}$, and $\widehat G_{ij}$ are parameters shared by all
rows of one aligned protein family.  Equations~\ref{eq:conditional-model} and
\ref{eq:reconstruction-loss} then define a sparse ternary pseudolikelihood model.  This branch is
the controlled setting for support and polarity recovery: synthetic ground truth can be specified,
MSA depth and phylogeny can be varied, and one position pair has a state comparable across all
sequences in the family.

Its limitation is scope.  The learned graph is tied to one alignment and does not by itself predict
relations in a new family or an unaligned sequence.  It is a statistical model of an ensemble, not
a replacement for a pretrained protein language model.

\subsection{Energy-compatible and general ESM-style mixers}

The strict associative version first projects source representations into memories
$x_j=K_h^{(\ell)}H_j^{(\ell)}$ and uses the same $x_j$ both to compute similarity and to return the
update in Eq.~\ref{eq:ternary-hopfield-update}.  If the memories, $\beta$, $\lambda_j$, and scaling
remain fixed during an inner recurrence, Proposition~\ref{prop:ternary-hopfield-descent} applies to
each target query.  This tied-memory head is the appropriate model for experiments that claim
fixed points, retrieval basins, or Lyapunov descent.  A learned output map may be applied after the
recurrence, but it is not itself part of the guaranteed descent dynamics.

The more expressive engineering version need not obey those constraints.

Let $H_i^{(\ell)}\in\R^d$ be the hidden state of residue $i$ at layer $\ell$.  For head $h$, an
ESM-style model \cite{rives2021} predicts three logits for every candidate edge,
\begin{equation}
r_{ij,h}^{z,(\ell)}
=f_{z,h}^{(\ell)}
\!\left(H_i^{(\ell)},H_j^{(\ell)},c_{ij}\right),
\qquad z\in\{-1,0,+1\},
\label{eq:dynamic-state-logits}
\end{equation}
where $c_{ij}$ may include relative position or declared side information.  Applying
Eq.~\ref{eq:state-softmax} on the state axis produces $m_{ij,h}^{(\ell)}$ and
$g_{ij,h}^{(\ell)}$.  A variance-controlled signed mixer is
\begin{align}
A_{i,h}^{(\ell)}
&=\sum_j m_{ij,h}^{(\ell)},\nonumber\\
M_{i,h}^{(\ell)}
&=\frac{\gamma_h}{\sqrt{\max\{1,A_{i,h}^{(\ell)}\}}}
\sum_j g_{ij,h}^{(\ell)}V_h^{(\ell)}H_j^{(\ell)},
\label{eq:dynamic-ternary-mixer}
\end{align}
followed by the usual output projection, residual update, normalization, and feed-forward block.
The normalization is over active mass rather than a probability simplex over sources.  A hard
version replaces the soft means by $z_{ij,h}^{(\ell)}$ and must use a sparse kernel to turn zero
states into measured compute savings.

Because Eq.~\ref{eq:dynamic-ternary-mixer} uses independent values $V_hH_j$, input-dependent
unconstrained logits, and active-mass normalization, it is generally not the gradient update of
Eq.~\ref{eq:ternary-hopfield-energy}.  Multi-head composition, residual connections, and changing
memories further remove the strict Lyapunov guarantee.  The tied head and the general head must
therefore be reported as two model variants rather than using an associative-memory proof to
justify the full ESM architecture.

The dynamic model is trained by
\begin{equation}
\mathcal L_{\mathrm{ESM}}
=\mathcal L_{\mathrm{MLM}}
+\lambda_0\sum_{\ell,h,i,j}m_{ij,h}^{(\ell)}
+\lambda_{\mathrm{cal}}\mathcal R_{\mathrm{soft/hard}}.
\label{eq:esm-objective}
\end{equation}
Unlike the family graph, its states are sequence- and layer-dependent.  Their scalar signs act on
shared value directions and are operational routing polarities, not intrinsic signs of unrestricted
amino-acid coupling tables.

There are two interpretation levels.  If Eq.~\ref{eq:dynamic-state-logits} observes the endpoint
representations, $z_{ij,h}^{(\ell)}$ is a dynamic routing decision.  A strict conditional
interaction claim requires logits constructed from endpoint-deleted context and an intervention
showing that gate parameters remain fixed while $X_i$ and $X_j$ vary.  Performance experiments may
use the ordinary dynamic version, but they must not silently inherit the stronger semantics of the
cavity version.

\subsection{Family teacher, sequence student}

The two branches can be connected without using an MSA at deployment.  Let $q_{ij}^{F}(z)$ and
$G_{ij}^{F}$ be detached state and field teachers fitted from a leave-query-out family alignment.
For an ESM student, add
\begin{equation}
\begin{split}
\mathcal L_{\mathrm{distill}}
={}&\sum_{i,j}\KL\!\left(
q_{ij}^{F}(z)\,\|\,
p_{ij}^{\theta}(z\mid X_{\setminus\{i,j\}})
\right)\\
&+\lambda_G\sum_{i,j}
\|G_{ij}^{\theta}-G_{ij}^{F}\|_{u_i,u_j}^2.
\end{split}
\label{eq:family-student-distillation}
\end{equation}
The leave-query-out construction prevents the target sequence from being copied into its own
teacher.  The field term requires the same categorical gauge on both sides.  At inference both
teacher terms are absent; the model receives one sequence.  This tests whether family-level
ternary organization is predictable from transferable sequence representations rather than tied
permanently to an MSA.

\section{Statistical-physics hypothesis}

For a paper aimed at PRL or PRX, the contribution cannot be the three-way softmax itself.  The
central testable hypothesis is a recovery transition controlled by sequence information and graph
complexity.  In synthetic ensembles, vary
\begin{equation}
\alpha=\frac{N_{\mathrm{eff}}}{L},\qquad
k,\qquad q,\qquad
\beta\rho,\qquad
\eta_{\mathrm{phy}},
\label{eq:control-variables}
\end{equation}
where $N_{\mathrm{eff}}$ is effective MSA depth, $k$ is true degree, $\beta\rho$ is interaction
strength, and $\eta_{\mathrm{phy}}$ controls phylogenetic distortion.  Measure support overlap,
polarity accuracy on oriented edges, operator error, and held-out reconstruction.

The associative branch has a separate phase diagram.  Sweep inverse temperature $\beta$, dead-zone
threshold $\lambda$, and memory load $N/d$; then measure energy descent, fixed-point type, retrieval
basin volume, and the transition between inactive, subset-composite, and single-memory retrieval.
Repeat after introducing phylogenetic correlations among memories.  This experiment tests a
hidden-state dynamical law and must not be conflated with the MSA sample-complexity variable
$N_{\mathrm{eff}}/L$ in Eq.~\ref{eq:control-variables}.

The PRL-level result would be a reproducible transition or scaling collapse, for example
\begin{equation}
\mathcal O_{\mathrm{support}}
\simeq
F\!\left(
\frac{N_{\mathrm{eff}}}{Lk(q-1)^2},
\beta\rho,
\eta_{\mathrm{phy}}
\right),
\label{eq:scaling-hypothesis}
\end{equation}
with a ternary model occupying a different accuracy--sparsity frontier from binary support,
continuous signed routing, and dense pair models.  Equation~\ref{eq:scaling-hypothesis} is a
hypothesis to derive and test, not a result of the present draft.

Because a continuous energy already contains negative, zero, and positive values, the strongest
null is not ordinary nonnegative attention.  It is a spike-and-slab signed-energy model
\begin{equation}
\kappa_{ij}=b_{ij}a_{ij},
\qquad b_{ij}\in\{0,1\},\quad a_{ij}\in\R,
\label{eq:binary-support-signed-energy}
\end{equation}
which has exact support and unrestricted polarity without a three-class variable.  Ternarity is
scientifically useful only if it improves calibrated mechanism recovery, optimization, or realized
compute relative to this baseline at matched capacity.

\section{Required experiments}

\begin{table*}[t]
\centering
\small
\begin{tabular}{p{0.17\textwidth}p{0.27\textwidth}p{0.28\textwidth}p{0.19\textwidth}}
\toprule
System & Intervention & Primary measurements & Essential baselines \\
\midrule
Sparse Potts ensembles & Sweep $N_{\mathrm{eff}}$, $L$, $k$, coupling strength, and alphabet &
Support/polarity/operator recovery; likelihood; phase boundary & Dense Potts, $\ell_1$ Potts,
binary support plus signed energy, NRI-style three-edge-type model \\
Phylogenetic synthetic MSAs & Evolve identical equilibrium models on trees of varying depth &
False-edge rate and transition shift & Sequence reweighting, phylogenetic nulls \\
Higher-order controls & Add pure three-body terms at fixed pair statistics & Pairwise reconstruction
floor and gate stability & Pairwise Potts, hypergraph extension \\
Real protein families & Held-out families across length and MSA depth & Masked likelihood, contacts,
DMS, double-mutant epistasis & Potts, Transformer, continuous signed and sparse nonnegative routing \\
ESM-style pretraining & Replace source-softmax heads by Eq.~\ref{eq:dynamic-ternary-mixer}; sweep
depth, heads, and active degree & MLM perplexity, transfer, gate calibration, and active-edge scaling
& Standard ESM, binary gates, sparsemax/entmax \cite{martins2016,peters2019}, continuous signed mixer \\
Family-to-ESM transfer & Train with and without Eq.~\ref{eq:family-student-distillation}; hold out
families & Teacher-state prediction, operator response, DMS, and contact transfer & No teacher,
shuffled-family teacher, endpoint-visible teacher \\
Ternary associative dynamics & Sweep $(\beta,\lambda,N/d)$ and iterate tied-memory updates; add
controlled memory correlations & Energy monotonicity, fixed-point classes, basin size, phase
boundaries & Standard modern Hopfield averaging, untied-value mixer, continuous signed retrieval \\
Execution study & Harden gates and vary target degree & FLOPs, peak memory, wall time, soft--hard gap
& Dense implementation at matched parameters \\
\bottomrule
\end{tabular}
\caption{Minimum validation program.  All comparisons require matched parameter or compute budgets,
multiple seeds, uncertainty intervals, and held-out protein families.}
\label{tab:experiments}
\end{table*}

The principal ablations are: remove the zero state; remove the sign; replace the edge-state
softmax by a source-position softmax; absorb sign into an unrestricted field; fix versus learn the
orientation; replace hard ternary states by continuous signed coefficients or
Eq.~\ref{eq:binary-support-signed-energy}; and compare the mean gate with hard sampling and the
proper latent mixture in Eq.~\ref{eq:latent-energy-mixture}.  These tests
separate gains caused by sparsity, polarity, optimization, quantization, and typed edge capacity.
The family and ESM branches must also be evaluated independently: success of a static family graph
does not imply transferable single-sequence routing, while an ESM likelihood gain does not imply
recovery of a stable family interaction graph.

Real-MSA evaluation must use conservation-matched and phylogeny-aware null ensembles.  Otherwise a
gate can reconstruct family entropy or clade structure and still be mislabeled as an epistatic
relation.  Contact prediction is useful but insufficient: it compresses a typed categorical field
to one scalar per position pair.  Double-mutant measurements or held-out conditional responses are
more direct tests of reconstructed interaction functions.

The APC hypothesis in Eq.~\ref{eq:apc-interpretation} requires a separate audit.  Measure
$m_{ij}$, $g_{ij}$, $\rho_{ij}$, and $\|G_{ij}^{\mathrm{eff}}\|$ before and after scalar-score
compression; fit endpoint propensities $a_i$ on training sequences; and compare $a_i a_j$ with the
empirical APC term on held-out sequences and families.  The hypothesis predicts that the leading
rank-one component is concentrated in activity or magnitude, that consistently oriented signed
scores have a weaker separable background, and that subtracting the fitted activity null enriches
pair-specific residuals for contacts or double-mutant epistasis.  Failure of these predictions
would reject the ternary activity explanation even if conventional APC still improves ranking.

\subsection{Go/no-go criteria}

The two-scale program should proceed sequentially rather than treating model scale as evidence.
First, on synthetic and family-level data, ternary states must either shift the recovery boundary,
improve calibrated support or mechanism recovery, or achieve the same likelihood with a lower
realized degree than Eq.~\ref{eq:binary-support-signed-energy}.  If none occurs, the ternary
variable is an unnecessary reparameterization and is not a PRL/PRX result.

Second, the ESM branch must preserve masked-token and downstream accuracy after hardening while
reducing measured execution, or show that family supervision transfers to held-out families beyond
shuffled-teacher and endpoint-visible controls.  Soft sparsity without skipped computation, or
teacher agreement without held-out biological improvement, does not pass this criterion.

Third, any claimed statistical-physics law must survive variation in $L$, $q$, family, phylogenetic
tree shape, and effective depth.  A threshold observed at one architecture size or one Pfam family
is a model-selection curve, not evidence of universality.

\section{Discussion}

The proposed model is best understood as a sparse signed categorical graphical model with a
neural, differentiable state selector.  It is not obtained by replacing the amino-acid output
softmax with the numbers $-1,0,+1$.  The output softmax defines conditional probabilities and hence
an effective-energy coordinate; the ternary softmax instead places a distribution over the state
of each candidate relation.  Neither softmax alone establishes a physical Hamiltonian.

The inactive state is potentially the most consequential component because it couples a
statistical decision to realizable sparse execution.  The sign can improve optimization or permit
a compact polarity code when value directions are oriented or otherwise restricted.  With
unrestricted categorical fields it is representationally redundant, so any claimed gain must be
localized by ablation.  This limitation is useful: it makes the hypothesis falsifiable and avoids
using a neuronal analogy as a substitute for an identified mechanism.

There are therefore two defensible semantics.  Under a static-energy interpretation, ternarity is
a quantized support controller layered on an otherwise continuous signed interaction, and the
binary-support signed-energy model in Eq.~\ref{eq:binary-support-signed-energy} is the decisive
baseline.  Under a latent-regime interpretation, $(p^-,p^0,p^+)$ represents uncertainty over three
mechanisms and must be propagated by categorical sampling or probability-level marginalization.
Using only $g=p^+-p^-$ is a useful relaxation but cannot by itself distinguish a physical mixture
from one weaker effective coupling.

A concise PRL manuscript should center on one derived and observed recovery transition, with the
ESM branch serving as a decisive universality or transfer test rather than a second unrelated
story.  A PRX manuscript could treat both scales symmetrically and would additionally require a
broader phase diagram, analytical sample-complexity theory, phylogenetic perturbations,
higher-order failure modes, multiple protein families, and a demonstrated compute frontier.  The
general reference-centered operator theory remains in the separate framework manuscript and should
be cited only for the definition of typed edge fields.

\section{Conclusion}

A ternary relation model can separate whether an edge is active from uncertainty about how an
oriented active edge acts.  It does not create attraction and repulsion that are absent from a
continuous signed energy, and its soft expected state is not exact marginalization over three
Hamiltonians.  At family scale it defines a sparse conditional reconstruction model for an MSA;
inside an ESM-style network it defines dynamic signed routing for single sequences.  Their
connection can be tested by leave-query-out distillation without requiring an alignment at
inference.  Its scientific value will be determined by calibrated recovery, transferable relation
states, or a new accuracy--compute frontier over binary-support signed-energy and continuous
baselines, not by the existence of three labels.

\begin{thebibliography}{99}

\bibitem{morcos2011}
F.~Morcos et al., ``Direct-coupling analysis of residue coevolution captures native contacts
across many protein families,'' \emph{PNAS} \textbf{108}, E1293--E1301 (2011).

\bibitem{ekeberg2013}
M.~Ekeberg et al., ``Improved contact prediction in proteins: Using pseudolikelihoods to infer
Potts models,'' \emph{Physical Review E} \textbf{87}, 012707 (2013).

\bibitem{dunn2008}
S.~D. Dunn, L.~M. Wahl, and G.~B. Gloor, ``Mutual information without the influence of phylogeny or
entropy dramatically improves residue contact prediction,'' \emph{Bioinformatics} \textbf{24},
333--340 (2008).

\bibitem{vaswani2017}
A.~Vaswani et al., ``Attention Is All You Need,'' in \emph{Advances in Neural Information
Processing Systems 30} (2017).

\bibitem{ramsauer2021}
H.~Ramsauer et al., ``Hopfield Networks Is All You Need,'' in \emph{Proceedings of ICLR} (2021).

\bibitem{rives2021}
A.~Rives et al., ``Biological Structure and Function Emerge from Scaling Unsupervised Learning to
250 Million Protein Sequences,'' \emph{PNAS} \textbf{118}, e2016239118 (2021).

\bibitem{rao2021}
R.~Rao, J.~Liu, R.~Verkuil, J.~Meier, J.~F. Canny, P.~Abbeel, T.~Sercu, and A.~Rives,
``MSA Transformer,'' in \emph{Proceedings of ICML}, PMLR \textbf{139}, 8844--8856 (2021).

\bibitem{dauparas2019}
J.~Dauparas, H.~Wang, A.~Swartz, P.~Koo, M.~Nitzan, and S.~Ovchinnikov,
``Unified Framework for Modeling Multivariate Distributions in Biological Sequences,''
ICML Workshop on Computational Biology, arXiv:1906.02598 (2019).

\bibitem{jang2017}
E.~Jang, S.~Gu, and B.~Poole, ``Categorical Reparameterization with Gumbel--Softmax,'' in
\emph{Proceedings of ICLR} (2017).

\bibitem{maddison2017}
C.~J. Maddison, A.~Mnih, and Y.~W. Teh, ``The Concrete Distribution: A Continuous Relaxation of
Discrete Random Variables,'' in \emph{Proceedings of ICLR} (2017).

\bibitem{li2016}
F.~Li, B.~Liu, X.~Wang, B.~Zhang, and J.~Yan, ``Ternary Weight Networks,''
arXiv:1605.04711 (2016).

\bibitem{martins2016}
A.~F. T. Martins and R.~F. Astudillo, ``From Softmax to Sparsemax: A Sparse Model of Attention and
Multi-Label Classification,'' in \emph{Proceedings of ICML}, PMLR \textbf{48}, 1614--1623 (2016).

\bibitem{peters2019}
B.~Peters, V.~Niculae, and A.~F. T. Martins, ``Sparse Sequence-to-Sequence Models,'' in
\emph{Proceedings of ACL}, 1504--1519 (2019).

\bibitem{kipf2018}
T.~N. Kipf, E.~Fetaya, K.-C. Wang, M.~Welling, and R.~S. Zemel, ``Neural Relational Inference for
Interacting Systems,'' in \emph{Proceedings of ICML}, PMLR \textbf{80}, 2688--2697 (2018).

\bibitem{ye2024}
T.~Ye, L.~Dong, Y.~Xia, Y.~Sun, Y.~Zhu, G.~Huang, and F.~Wei, ``Differential Transformer,''
arXiv:2410.05258 (2024).

\bibitem{chen2026}
J.~Chen, G.~Li, J.~E. Hopcroft, and K.~He, ``SignGT: Signed Attention-Based Graph Transformer for
Graph Representation Learning,'' \emph{Knowledge and Information Systems} \textbf{68}, 221
(2026).

\end{thebibliography}

\end{document}
